\chapter{Introduzione}
La Ricerca Operativa è la disciplina che studia e mette a punto metodologie ed algoritmi per la risoluzione di problemi di ottimizzazione
\begin{definition}{Problema Ottimizzazione}
    Un problema di ottimizzazione è un problema in cui occorre prendere decisioni sull'uso di risorse disponibili solo in quantità limitate,
    in modo tale da rispettare un dato insieme di vincoli, e da massimizzare il beneficio ottenibile dall'uso delle risorse \textbf{o minimizzare
    il costo totale che ne deriva dall'utilizzo di queste risorse.}
\end{definition}
\begin{example}{Esempio di problema di ottimizzazione}
    Supponiamo che un'azienda produca tavoli e sedie:
    \begin{itemize}
        \item Ogni tavolo richiede 4 ore di lavoro e 2 kg di legno.
        \item Ogni sedia richiede 2 ore di lavoro e 1 kg di legno.
        \item L'azienda dispone di 40 ore di lavoro e 16 kg di legno.
        \item Guadagna 50 € per ogni tavolo e 20 € per ogni sedia.
    \end{itemize}
    Indichiamo con
    \begin{itemize}
        \item x = numero di tavoli prodotti.
        \item y = numero di sedie prodotte.
    \end{itemize}
    Vogliamo massimizzare il profitto
    \begin{center}
        $\max_{x,y} (50x + 20y)$


        soggetto ai vincoli
        $$
        \begin{cases}
        4x + 2y \le 40 & \text{(ore di lavoro)} \\
        2x + y \le 16 & \text{(legno)} \\
        x, y \ge 0
        \end{cases}
        $$
    \end{center}
\end{example}
Quindi un processo decisionale alla base dell'analisi e risoluzione di un problema
di ottimizzazione puo essere scomposto nelle seguenti fasi:
\begin{itemize}
    \item individuazione del problema e delle sue caratteristiche;
    \item costruzione di un modello logico-matematico che lo rappresenti (formulazione degli obiettivi e scelta della variabili di interesse, dette variabili
di decisione, nell'esempio di prima sono x ed y);
    \item determinazione di una o piu soluzioni;
    \item analisi dei risultati ottenuti.
\end{itemize}
\newpage
\section{Classificazione dei problemi di ottimizzazione}
Per definire matematicamente un problema di ottimo, quello che si fa è definire:
\begin{definition}{Insieme ammissibile}
    \begin{center}
        P = \{\text{tutte le soluzioni che rispettano i vincoli del problema}\}
    \end{center}
    P viene chiamato insieme ammissibile e ogni elemento x $\in$ P è \textbf{una soluzione ammissibile}
\end{definition}
\begin{example}{Esempio di Insieme ammissibile}
    Nell'esempio dei tavoli e delle sedie avevamo:
    \begin{center}
        $$
        \begin{cases}
        4x + 2y \le 40 & \text{(ore di lavoro)} \\
        2x + y \le 16 & \text{(legno)} \\
        x, y \ge 0
        \end{cases}
        $$
    \end{center}
    Quindi l'insieme ammissibile è 
    \begin{center}
        $$
        P = \left\{ (x,y) \in \mathbb{R}^2 : 
        \begin{array}{l}
        4x + 2y \le 40 \\
        2x + y \le 16 \\
        x \ge 0, y \ge 0
        \end{array} 
        \right\}
        $$
    \end{center}
\end{example}
\begin{definition}{Funzione Obiettivo}
    Sull'insieme $P$, viene definita una funzione:
    \begin{center}
        $c : P \to \mathbb{R}$
    \end{center}
    Questa funzione associa ad ogni soluzione ammissibile un valore reale:
    \begin{center}
           $x \in P \to c(x) \in \mathbb{R}$
    \end{center}
    Il valore \(c(x)\) rappresenta, a seconda del problema:
    \begin{itemize}
        \item \textbf{un profitto/beneficio}, che vogliamo massimizzare.
        \item \textbf{un costo}, vogliamo minimizzarlo.
    \end{itemize}
\end{definition}
\begin{definition}{Problema di Ottimizzazione}
    Quindi un problema di ottimizzazione può essere scritto semplicemente come:
    \begin{center}
        $\max_{x \in P} c(x)$
    \end{center}
    oppure
    \begin{center}
        $\min_{x \in P} c(x)$
    \end{center}
\end{definition}
\begin{example}{Esempio di funzione obiettivo}
    Un esempio di funzione obiettivo definita sull'insieme $P$ che abbiamo definito prima è:
    \begin{center}
        c(x,y) = 50x + 20y
    \end{center}
    Poichè vogliamo massimizzare il profitto, ho:
    \begin{center}
         $\max_{(x,y) \in P} 50x + 20y$
    \end{center}
\end{example}
\begin{definition}{Soluzione Ottima}
    Una soluzione ottima sarà un x* $\in P$ tale che, nel caso di massimizzazione,
    \begin{center}
        c(x*) $\geq$ c(x) $\forall x \in P$ 
    \end{center}
\end{definition}
In alcuni casi, invece di richiedere una soluzione ottima di un problema,
pu`o essere richiesto di individuarne una qualsiasi soluzione ammissibile, ossia
di determinare se l'insieme delle soluzioni ammissibili sia vuoto o meno. In
casi del genere, si parla di problema decisionale o problema di esistenza.
Ad un problema decisionale è naturalmente associato il cosiddetto problema
di certificato, che consiste nello stabilire se un dato x appartenga o meno a $P$
\section{Classificazione dei metodi risolutivi}
Gli algoritmi utilizzati per risolvere i problemi di ottimizzazione si suddivi
dono in tre categorie principali:
\begin{itemize}
    \item \textbf{Metodi esatti} :  essi risolvono il problema all'ottimo. In caso di prob
lemi computazionalmente intrattabili, essi richiedono complessit`a com
putazionale esponenziale o comunque superpolinomiale.
    \item \textbf{Metodi di approssimazione} : essi trovano una soluzione subottima,
garantendo un indice della qualit`a della soluzione approssimata indi
viduata. La loro applicazione risulta utile in caso di problemi computazionalmente intrattabili. è chiaro che in tal caso sarebbe auspicabile che il tempo di computazione da essi richiesto sia polinomiale.
    \item \textbf{Metodi euristici} :  essi trovano una soluzione subottima senza alcuna
garanzia circa la qualit`a della soluzione individuata, per cui `e ovvio che
ha senso la messa a punto di euristiche solo polinomiali. L'applicazione
di tali metodi `e particolarmente indicata per risolvere problemi del
mondo reale di grandi dimensioni e/o problemi talmente difficili dal
punto di vista computazionale che per essi non esiste neanche un al
goritmo capace di trovare una soluzione approssimata in tempo poli
nomiale.


\end{itemize}